\chapter{Apr.~20 --- Geometric Aspects}

\section{Local Systems}

\begin{remark}
  Let $X$ be a real manifold, and
  let $\Omega^i(X)$ denote the space of
  differential $i$-forms on $X$.
  Let
  \[
    \Omega^i(X, E)
    = E \otimes \Omega^i(X)
  \]
  denote the space of differential
  forms with values in a vector bundle $E$.
\end{remark}

\begin{definition}
  A \emph{connection} on $E$ is a map
  $\nabla : \Gamma(U, E) \to \Omega^1(U, E)$
  satisfying
  \[
    \nabla(fs) = f \nabla s + s \otimes df
  \]
  for a function $f$ and a section $s$.
\end{definition}

\begin{remark}
  We can define
  $\nabla_\xi : E \to E$
  by $s \mapsto \langle \nabla s, \xi \rangle$
  with
  \[
    \nabla_\xi(fs)
    = f \nabla_\xi s + (\partial_\xi f) s.
  \]
  In local coordinates, we can write
  \[
    \nabla_i
    = \nabla_{\partial / \partial x_i}
    = \partial_{x_i} + A_i(x).
  \]
  If $U_{\alpha, \beta}$ are the transition
  functions for $E$, then
  \[
    U_{\alpha, \beta}
    (d + A_\alpha) U_{\alpha, \beta}^{-1}
    = \nabla_\beta = d + A_\beta.
  \]
  This implies
  $A_\beta = U_{\alpha, \beta} A_\alpha U_{\alpha, \beta}^{-1} + U_{\alpha, \beta} d U_{\alpha, \beta}^{-1}$,
  and such a map is
  called a \emph{gauge transformation}.
  Note that we get a \emph{curvature 2-form}
  by defining $F = dA + A \wedge A$,
  and we have
  \[
    F(\xi, \eta)
    = [\nabla_{\xi}, \nabla_\eta] - \nabla_{[\xi, \eta]}.
  \]
  The components of the 2-form $F$
  are given by
  \[
    [\partial_i + A_i, \partial_j + A_j]
    = \partial_i A_j - \partial_j A_i + [A_i, A_j]
    = F_{i, j}.
  \]
\end{remark}

\begin{definition}
  A \emph{flat} connection is a connection
  with $0$ curvature, i.e.
  \[
    [\nabla_\xi, \nabla_\eta] = \nabla_{[\xi, \eta]},
  \]
  or equivalently, $F_{i, j} = 0$.
  A \emph{local system} is a vector bundle
  which admits a flat connection.
\end{definition}

\begin{remark}
  A connection $\nabla$ gives rise to
  \emph{parallel transport}. Given a
  curve $\gamma : [0, 1] \to X$ and
  $v \in E_{\gamma(0)}$ (the fiber of
  $E$ at $\gamma(0)$), there is a unique
  section $s$ of $E$ along $\gamma$ such that
  $\nabla_{\dot{\gamma}} s = 0$ and
  $s(\gamma(0)) = v$.
  There is a \emph{holonomy operator}
  $M_{\gamma} : E_{\gamma(0)} \to E_{\gamma(1)}$
  given by $s(\gamma(0)) \mapsto s(\gamma(1))$.
\end{remark}

\begin{theorem}
  If $\nabla$ is a flat connection, then
  $M_\gamma$ depends only on the homotopy
  class of $\gamma$.
\end{theorem}

\begin{corollary}\label{cor:local_system}
  The map
  $\gamma \mapsto M_\gamma$
  (with fixed basepoint) gives a
  representation of the fundamental group
  in $E_{x_0}$.
\end{corollary}

\begin{remark}
  Building on Corollary \ref{cor:local_system}, there is in fact
  a bijection
  between local systems and finite-dimensional
  representations of $\pi_1(X, x_0)$.
\end{remark}

\begin{definition}
  Let $U \subseteq X$ and $s$ a section
  of $E$ over $U$. We say that $s$ is
  \emph{flat} if
  $\nabla_\xi s = 0$ for all
  $\xi \in U$. We denote the space of
  flat sections over $U$ by
  $\Gamma_f(U, E)$.
\end{definition}

\begin{lemma}
  Let $U \subseteq X$ be open and
  simply connected. Then for any $p \in U$,
  there is an isomorphism
  \begin{align*}
    \Gamma_f(U, E)
    &\longrightarrow E_p \\
    s &\longmapsto s(p).
  \end{align*}
  In other words, the sheaf of flat sections is
  locally constant.
\end{lemma}

\begin{remark}
  Let $E, E'$ be two vector bundles
  over $X$ with connections
  $\nabla, \nabla'$. Then
  \[
    \Delta(s \otimes s')
    = (\nabla s \otimes s')
    + (s \otimes \nabla' s')
  \]
  defines a connection on $E \otimes E'$.
  Similarly, if $E^*$ is dual to $E$
  with pairing $\langle s, s^* \rangle$, then
  \[
    \langle \nabla_\xi s, s^* \rangle
    + \langle s, \nabla_\xi s^* \rangle
    = \partial_\xi \langle s, s^* \rangle
  \]
  defines a connection on $E^*$.
  One can also define a pullback bundle
  $f^* E$ and
  corresponding connection, given a map
  $f : Y \to X$ and vector bundle $E$ on $X$.
\end{remark}

\section{De Rham Cohomology}

\begin{remark}
  We have a sequence
  \[
    \begin{tikzcd}
      0 \ar[r] & E \ar[r, "\nabla"] & \Omega^1(X, E) \ar[r, "\nabla"] & \Omega^2(X, E) \ar[r, "\nabla"] & \cdots \ar[r, "\nabla"]
               & \Omega^n(X, E) \ar[r] & 0
    \end{tikzcd}
  \]
  where $\nabla(s \otimes \omega) = \nabla(s) \wedge \omega + s \otimes d\omega$.
  However, this is not a complex in general.
\end{remark}

\begin{prop}
  If $\nabla$ is a flat connection, then
  $\nabla^2 = 0$ and
  $\Omega^\bullet(X, E)$ is a complex.
\end{prop}

\begin{remark}
  If $\nabla$ is flat, then we
  can consider the de Rham cohomology
  $H^\bullet_{\mathrm{dR}}(X, E)$
  with coefficients in a local system
  (the cohomology of the
  complex $\Omega^\bullet(X, E)$).
\end{remark}

\begin{remark}
  Let $\sigma : \Sigma_k \to X$ be a
  singular $k$-simplex. Consider
  $\sum_j \varepsilon_j \sigma_j$, where
  $\varepsilon_j$ is a flat section 
  of $\sigma^* E$ over $\Sigma_k$. Then
  we can define a boundary map $\partial$ by
  \[
    \partial (\varepsilon_j \sigma_j)
    = \sum_j \sum_\ell \varepsilon_j^\ell \sigma_j^\ell (-1)^\ell,
  \]
  where $\varepsilon_j^\ell$ is the restriction
  of the section to the $\ell$th face and
  $\sigma_j^k$ is the $\ell$th face of
  $\sigma_j$. This gives rise to the
  \emph{singular homology}
  $H_\bullet(X, E)$
  and \emph{singular cohomology}
  $H^\bullet(X, E)$.
\end{remark}

\begin{remark}
  Let $E$ be a local system with
  fiber $V$. Let $\omega \in \Omega^k(X, E)$
  and
  $\sigma : \Sigma_k \to X$, so
  $\sigma^* \omega$ is a differential form
  on $\Sigma_k$ with coefficients in the
  local system $\sigma^* E$. If
  \[
    \omega \in \Omega^k(\Sigma_k, \sigma^* E)
    = \Omega^k(\Sigma_k) \otimes V
  \]
  and $C = \varepsilon \sigma$ (where
  $\varepsilon$ is the coefficient section
  of the dual bundle), then we can define
  \[
    \int_C \omega
    = \langle \varepsilon, \int_{\Sigma_k} \sigma^* \omega \rangle.
  \]
\end{remark}

\begin{theorem}[Stokes theorem]
  We have
  \[
    \int_C d\omega = \int_{\partial C} \omega,
  \]
  and integration defines a pairing
  $B : H^\bullet(X, E^*) \otimes H^\bullet_{\mathrm{dR}}(X, E) \to \C$.
\end{theorem}

\begin{theorem}[De Rham theorem]
  If $X$ is a smooth real manifold and
  $B$ is nondegenerate, then
  $H^\bullet(X, E^*) \cong H^\bullet_{\mathrm{dR}}(X, E)$.
\end{theorem}


\begin{remark}
  Consider an \emph{affine complex manifold},
  i.e. a complex manifold $X$ which is
  isomorphic to a closed subvariety
  in $\C^N$. We can
  write $X = \C^N \setminus D$ for
  some subvariety $D \subseteq \C^N$ of
  codimension $1$.
\end{remark}

\begin{theorem}[Complex de Rham theorem]
  Let $X$ be an affine complex manifold
  and $H^{\bullet}_{\mathrm{dR}}(X, E)$ be the
  cohomology of the global holomorphic
  forms. Then $B$ is nondegenerate and
  $H^\bullet(X, E^*) \cong H^\bullet_{\mathrm{dR}}(X, E)$.
\end{theorem}

\begin{corollary}
  For a complex affine manifold of
  complex dimension $n$, we have that
  $H_i(X, E)$ and $H^i(X, E)$
  are nonzero only for $i \le n$.
\end{corollary}

\section{Hyperplane Arrangements}

\begin{remark}
  Consider a finite set $S$ of codimension
  $1$ affine hyperplanes
  in $\C^n$ (a \emph{hyperplane arrangement}), and let $X$ be the complement
  of its union. Let $H \in S$ and
  $\ell_H$ the corresponding linear
  functional (i.e. $H$ is the set of
  solutions to $\ell_H(\vec{z}) = 0$).
  Let $\vec{\lambda} = \{\lambda_H\}_{H \in S}$
  for some constants $\lambda_H \in \C$
  and define
  \[
    \psi(\vec{z})
    = \prod_{H \in S} \ell_H(\vec{z})^{\lambda_H}.
  \]
\end{remark}

\begin{definition}
  For any $S$ and $\vec{\lambda}$,
  denote by $\mathcal{L}_{\vec{\lambda}}$ the
  $1$-dimensional local system
  over $X$ with local sections
  $\psi(\vec{z}) f(\vec{z})$
  (where $f$ is a holomorphic function on
  $X$) and connection
  $\nabla_i = \partial_i$. Equivalently,
  we can define $\mathcal{L}_{\vec{\lambda}}$
  to be the trivial line bundle with
  connection
  \[
    \nabla_i f
    = \partial_i f + f \partial_i(\log \psi)
    = \partial_i f + f \psi^{-1} \partial_i \psi.
  \]
\end{definition}

\begin{definition}
  Define the \emph{Orlik-Solomon algebra}
  $\mathrm{OS}^\bullet(X)$ to be the
  graded exterior algebra generated
  by $d(\log \ell_H(\vec{z}))$,
  which can be described by generators
  and relations.
\end{definition}

\begin{remark}
  For $f \in \mathrm{OS}^\bullet(X)$, we call
  $\psi f$ a \emph{hypergeometric form},
  and denote the space of hypergeometric
  forms by
  $\Omega^\bullet_{\mathrm{hg}}(X, \mathcal{L}_{\vec{\lambda}}) \subseteq \Omega^\bullet(X, \mathcal{L}_{\vec{\lambda}})$.
\end{remark}

\begin{theorem}\label{thm:hypergeometric}
  For almost all $\vec{\lambda}$
  (including $\vec{\lambda} = 0$), the map
  \[
    \Omega^\bullet_{\mathrm{hg}}(X, \mathcal{L}_{\vec{\lambda}})
    \longrightarrow \Omega^\bullet(X, \mathcal{L}_{\vec{\lambda}})
  \]
  is a quasi-isomorphism.
\end{theorem}

\begin{example}\label{ex:hypergeometric}
  Let $X = \C \setminus \{0\}$ and
  $\psi(z) = z^\lambda$. Then
  \[
    z^\lambda d(\log z)
    = z^\lambda \frac{dz}{z}
  \]
  is a hypergeometric 1-form, and
  $z^\lambda$ is a $0$-form.
  Note that
  \[
    d(z^\lambda)
    = \lambda z^{\lambda} \frac{dz}{z}.
  \]
  We have $H^1 = \C$ for
  $\lambda = 0$ and $H^1 = 0$ for
  almost all $\lambda \ne 0$.
\end{example}

\begin{exercise}
  Show that in
  Example \ref{ex:hypergeometric}, we have
  $H^1 = \C$ if $\lambda \in \Z$
  and $H^1 = 0$ if $\lambda \notin \Z$.
  So we see that Theorem
  \ref{thm:hypergeometric}
  holds for
  $\lambda \notin \Z \setminus \{0\}$.
\end{exercise}

\begin{remark}
  One can consider
  discriminantal arrangements on configuration spaces.
  Let $\vec{z} = (z_1, \dots, z_N) \in \C^N$
  such that $z_i \ne z_j$. For
  $m \ge 0$, we have a map
  $X_{m + N} \to X_N$ given by forgetting
  all but the first $N$ components.
  This map is a fiber bundle.
  Denote the fiber over $\vec{z} \in X_N$
  by $Y_{\vec{z}, m}$, i.e.
  \[
    Y_{\vec{z}, m}
    = \{\vec{t} = (t_1, \dots, t_m) \in \C^m :
    t_i \ne t_j,\, t_i \ne z_p\}.
  \]
  Moreover, there is a natural action of
  $S_N$ on $X_N$ and of $S_m$ on
  $Y_{\vec{z}, m}$ by permuting
  coordinates. Let
  \[
    Q = \{
      q_{p, n}, r_{p, j} : 1 \le p < n \le m,\, 1 \le j \le N
    \}.
  \]
  Then we can define
  \[
    \psi_{\vec{z}}(\vec{t})
    = \prod_{p < n} (t_p - t_n)^{q_{p, n}}
    \prod_{p, j} (t_p - z_j)^{r_{p, j}}.
  \]
  If $q_{p, n}, r_{p, j}$ are symmetric
  with respect to $S_m$, then
  this defines a local system
  $\overline{\mathcal{L}}_Q$
  on
  $\overline{Y}_{\vec{z}, m}
    = Y_{\vec{z}, m} / S_m$.
  Moreover, we have
  $H^\bullet(\overline{Y}_{\vec{z}, m}, \overline{\mathcal{L}}_Q) = (H^\bullet(Y_{\vec{z}, m}, \mathcal{L}_Q))^{S_m}$.
\end{remark}

\begin{remark}
  Consider the KZ equations for
  $\mathfrak{sl}_2$. Let
  $\vec{z} = (z_1, \dots, z_N)$,
  $\vec{t} = (t_1, \dots, t_m)$,
  and $\kappa \in \C^\times$.
  Let $\mathcal{L} = \mathcal{L}(\vec{\mu}, \kappa)$.
  Then recall that we have
  \[
    \psi_{\vec{m}}(\vec{z}, \vec{t})
    = \prod_{i < j} (z_i - z_j)^{\mu_i \mu_j / 2\kappa}
    \prod_{p, j} (t_p - z_j)^{-\mu_j / \kappa}
    \prod_{p < n} (t_p - t_n)^{2 / \kappa}.
  \]
\end{remark}
